% HD-R02--R03. Root construction, 21 November 2026.
% Whole-proof audit and HD-G18 row certificate are recorded separately.
\subsection{HD-R02: an explicit integral detector on the nonsplit Gaussian source}
\label{sec:HD-R02}

Use the actual group, curve and primitive factor of HD-S22--S23, and put
\[
 \begin{gathered}
 G=1393459200,\quad g=3077222401,\quad D=174182400,\\
 N=10^{755},\qquad L=2gN,\qquad k=gN+3.
 \end{gathered}
\]
Here $g$ denotes the genus of the Galois curve, while $G$ is its group
order. Define the following explicit positive integers:
\begin{align}
 S_{\rm ord}&=G^{2g}(g!)^2,\nonumber\\
 S_{\rm prod}&=(156341088!)^2\,12^{312682176}
                         C_{12}(609948,45),\nonumber\\
 S_0&=S_{\rm ord}+S_{\rm prod},\qquad
 \lambda=8(g+G)=35765452808,\nonumber\\
 \mathcal B&=\lambda^k\binom{k}{3}(gN)!S_0^N\,720D^6.
                                                        \tag{R02.1}
\end{align}
The integer $C_{12}$ is the closed factorial expression (N15.3).
Every primitive sixfold $A$ in the E8 source component has an integral
effective three-cycle with nonzero Weil projection and degree at most
$\mathcal B$ for the restricted Jacobian theta polarization.
All the factors in (R02.1) are displayed integers; this is not a
minimum over an unspecified collection of cycles.

\paragraph{1. The exact supplier identity and its sizes.}
HD-G18 identifies the selected $98$ rows of the candidate integer
identity with actual subgroup characters of the actual covering group.
The fresh integer calculation gives
\[
             \sum_{j=1}^{98}a_jv_j=e_\rho,
             \qquad \sum_j|a_j|<10^{755}=N.           \tag{R02.2}
\]
The file \texttt{witness\_identity.json} in \texttt{work/root/e8/}
retains the literal coefficients, rows and exact product.
The subgroup realizations and their character maps are separate
certificates in HD-G18.
The source identity is credited to B.~Severn's CC~BY~4.0 candidate
archive; the integer identity and every selected subgroup realization
are verified here.

There are $95$ ordinary admissible cyclic rows, including the
Sylow-three normalizer row, and three product rows. An ordinary row
uses $h_j\le g$ curve factors and has signed degree mass at most
$S_{\rm ord}$ by HD-S25. For the three product rows, the freshly
constructed positive slot lists give
\[
\begin{array}{c|r|r|r|r}
\text{row}&R&H&S&\text{genus of the quotient curve}\\\hline
3231&2&19784&6748&3259\\
3244&45&156341088&609858&1730347\\
3269&4&42739160&4816&5341793
\end{array}                                            \tag{R02.3}
\]
Here $S=\sum(s_j-2)$, so $H=2g_XR+S$ in each row.
All slots have $s_j\ge2$ and character order dividing $12$;
the four simultaneous conjugates have equal Hodge numbers.
HD-R01 therefore supplies effective middle cycles of degree at most
$(H!)^2 12^{2H}C_{12}(S,R)$.
The formula (N15.3) increases when its even first argument or its
second argument increases. Thus the deliberately larger first
argument $609948$ in (R02.1) bounds all three rows. Their $H$ is at
most $156341088<g$, so every selected supplier uses at most $g$
curve factors and has mass at most $S_0$.

\paragraph{2. The determinant character map on the full centralizer.}
Use the left-deck Jacobian action $a_h=L_{h,*}$ and the cohomological
action $T_h=(L_h^{-1})^*$ fixed in HD-G18. Write
$V=H^1(C,\mathbb C)$ with this $T$-action and
\[
              U=\operatorname{Cent}_{\operatorname{Sp}(V)}(G).
\]
For an irreducible complex $G$-representation $\tau$, put
$M_\tau=\operatorname{Hom}_G(\tau,V)$.
The cup form pairs $M_\tau$ with $M_{\bar\tau}$ perfectly.
Every real irreducible representation of this particular $G$ has a
symmetric invariant form: HD-G18 verifies all $112$ Frobenius--Schur
indicators in the Weyl group, and the real characters of the central
product extend these real representations with central scalar $\pm1$.
Their multiplicity spaces are symplectic. Consequently
\[
 U\simeq\prod_{\tau\text{ real}}\operatorname{Sp}(M_\tau)
       \ \times\!
       \prod_{\{\tau,\bar\tau\}\text{ nonreal}}
                         \operatorname{GL}(M_\tau).   \tag{R02.4}
\]
Zero multiplicities are omitted. This describes the actual full
centralizer and its characters; it is not a replacement by a smaller
connected subgroup.

For a linear character $\chi$ of $H\le G$, Frobenius reciprocity is
the actual evaluation isomorphism
\[
 \operatorname{Hom}_G(\operatorname{Ind}_H^G\chi,V)
          \longrightarrow\operatorname{Hom}_H(\chi,V).
\]
Taking its top exterior power shows that the exponent of
$\det M_\tau$ is
$\langle\operatorname{Res}_H\tau,\chi\rangle-
 \langle\operatorname{Res}_H\bar\tau,\chi\rangle$.
These are precisely the $45$ coordinates used in HD-G18.
The symplectic factors have determinant one. Products of slots add
these exponents. Thus (R02.2) is exactly an identity of characters
of $U$, with no discarded sign character.

\paragraph{3. A pure cohomological tensor paired with actual cycles.}
For each row choose its signed integral middle cycle $Y_j$ from
HD-S25 or HD-R01, on the specified product $P_j$ of $h_j$ quotient
curves. Its mass is at most $S_0$, and it evaluates nontrivially on
the inverse determinant tensor of that row. The curve maps used
here are the actual maps
\[
 C/\ker\chi\longrightarrow J(C/\ker\chi)
                      \longrightarrow J(C),          \tag{R02.5}
\]
where the second arrow is pullback on Jacobians.
For a slot with $h$ curve factors, first pull its determinant class
back under addition $J(C)^h\to J(C)$. Restriction to the curve
product is the alternating determinant tensor, up to the nonzero
factor $h!$: two one-forms with the same order-at-least-three
character on one curve have product in $H^2_{\chi^2}=0$.
For $p:C\to C'=C/\ker\chi$, write $f=p_J^*\circ
\operatorname{AJ}_{C'}$ for (R02.5). On one-forms its pullback
is the trace $f^*=p_*$, and $p_*p_{\rm curve}^*=[\deg p]$.
Identify $H^1(C')_\chi$ inside $H^1(C)$ by the actual curve
pullback $p_{\rm curve}^*$. On this identified summand $f^*$ is
therefore the nonzero scalar $\deg p$, and on its determinant
it is the corresponding positive power. Frobenius reciprocity
identifies this summand with $\operatorname{Hom}_H(\chi,V)$
through evaluation. These are the exact maps giving the nonzero
evaluations.

Use $|a_j|$ copies of $P_j,Y_j$. Choose the inverse determinant
line for $a_j>0$ and the determinant line for $a_j<0$.
Rationality of the cycle gives nonzero evaluation for the conjugate
line as well. Denote their exterior products by $P,Y$, and let
$\ell=\sum_j|a_j|h_j$. Then $\ell$ is even, $\ell<gN\le L$,
and $Y$ has dimension $\ell/2$.
Let $\zeta\in H^\ell(J(C)^\ell,\mathbb C)$ be the product of the
chosen, addition-polarized determinant tensors. The preceding
evaluation gives
\[
                \int_Y\zeta\ne0,
       \qquad \operatorname{char}_U(\zeta)
                              =(\det M_\rho)^{-1}.    \tag{R02.6}
\]
Let $x\in H^6(A,\mathbb C)$ generate the corresponding Weil line.
The pure tensor $\Gamma=\zeta\boxtimes x$ is $U$-invariant and
its action on $Y$ is the nonzero scalar
$(\int_Y\zeta)x$.
This equality pairs the specified tensor with the actual cycle;
other character components of the cycle cannot alter it.

\paragraph{4. Explicit divisor products span the transfer tensor.}
Let $j:A\hookrightarrow J(C)$ be the inclusion and
$q=De:J(C)\to A$ the integral primitive projection. Their composite
is $qj=[D]_A$. Pull $\Gamma$ back by $1\times q$ to
$J(C)^{\ell+1}$. Its restriction by $1\times j$ is $D^6\Gamma$.
On cohomology $a_h^*=T_{h^{-1}}$, so the same finite set of
group correspondences spans the $T$-algebra. All these maps commute
with $U$.

The invariant theory of (R02.4) makes every invariant in its exterior
algebra a sum of products of invariant alternating forms of degree
two. The precise published input is Milne, \emph{Lefschetz classes
on abelian varieties}, Duke Math. J.~96 (1999), 639--675,
Lemma~3.5, Proposition~3.6 and Remark~3.7, author-version pp.~15--18.
Here is the additional degree-two identification used for the bound.
Let $B$ be the image of $\mathbb C[G]$ in
$\operatorname{End}_{\mathbb C}(V)$, the complex scalar extension
of the actual rational group-correspondence algebra.
Its centralizer is semisimple and stable under the cup adjoint.
Milne's Lemma~3.5 says that its unitary elements generate it as an
algebra. Taking the double centralizer therefore gives
$\operatorname{End}_U(V)=B$.
An invariant bilinear pairing is consequently
$(v,w)\mapsto\langle v,bw\rangle$ for $b\in B$.
Since $B$ is spanned by the actual elements of $G$, these forms
are spanned by the cross terms of theta pulled back under
$x_i+g_0x_j$, $g_0\in G$, together with the individual theta forms.
For repeated indices this gives the required alternating part by
the same pullback identity. This proves the spanning assertion for
the following finite list of actual divisors:
\[
        x_i^*(2\Theta_C),\qquad
        (x_i+g_0x_j)^*(2\Theta_C),\quad g_0\in G.    \tag{R02.7}
\]
The line bundle $2\Theta_C$ is globally generated. Explicitly,
choose a translate $a$ so that a specified point lies in neither
$\Theta_C+a$ nor $\Theta_C-a$; their sum is linearly equivalent
to $2\Theta_C$ by the theorem of the square.

Since $\Gamma$ has degree $\ell+6$, its pullback is a complex linear
combination of products of exactly $c=\ell/2+3$ divisors from
(R02.7). Restrict to $P\times A$ through the individual maps
(R02.5) and $j$. The resulting monomial contractions on $Y$ cannot
all have zero Weil projection: their linear combination is the
nonzero class $D^6(\int_Y\zeta)x$.
No denominator in this linear combination is inserted into any
cycle. One monomial suffices for nonvanishing.

\paragraph{5. Integral effectiveness and the numerical cap.}
Write $\eta$ for the sum of degree-one point classes on the
$\ell$ curve factors. The external-product degree formula gives
\[
 \operatorname{mass}_\eta(Y)
 \le (\ell/2)!\prod_j\operatorname{mass}(Y_j)^{|a_j|}
 \le(gN)!S_0^N.                                      \tag{R02.8}
\]
For a quotient $p:C\to C'$ of degree $t$, the map in (R02.5)
has theta degree $t g(C')\le g-1+G$ by Riemann--Hurwitz.
The invariant positive theta form gives
$\theta(v+w)\le2\theta(v)+2\theta(w)$.
Thus every restricted divisor from (R02.7), including repeated
indices, is nef-bounded by
$\lambda(\eta+\theta_A)$; this is the exact estimate of HD-S25.
Choose general sections successively so that they intersect all
positive and negative components properly, and push their integral
effective intersections to $A$.
The projection formula gives the mass bound
\[
 \lambda^c\binom{c}{3}\operatorname{mass}_\eta(Y)
                           \int_A\theta_A^6.
\]
The binomial term is exact: in
$(\eta+\theta_A)^c\theta_A^3$ only the term of curve degree
$\ell/2$ and abelian degree $6$ survives.
HD-S23 gives $\int_A\theta_A^6=720\Delta_A\le720D^6$.
Since $c\le k$, (R02.8) proves the bound (R02.1).
Choose the monomial with nonzero Weil projection from section~4,
then an effective component of one of its two signed parts with
nonzero projection. Its degree is at most the displayed mass.
This is the required effective integral detector.

\subsection{HD-R03: target-lattice and division bounds for determinant class $[-3]$}
For a specified Gaussian Weil lattice of signature $(3,3)$ and
determinant class $[-3]$, let $M$ be the terminating, explicitly
computed denominator in HD-S24. Its actual isogeny has
$\phi^*\theta_t=M^2\theta_A$. Consequently on a Zariski-dense
set of this Weil family,
\begin{equation}
          \delta_W\le M^6\mathcal B.                 \tag{R03.1}
\end{equation}
Indeed the smooth rank-nine point in HD-G17 gives an analytic open
period image. Transport by the fixed rational isometry of HD-S24
preserves that openness, and the displayed isogeny sends every
source detector to an integral effective target detector with
degree multiplied by $M^6$. An analytic open subset of the
connected Weil family is Zariski dense. The supplied reduction
R1--R2 and Lemma Dom then give algebraicity and the same uniform
bounded-division consequence on the whole fixed family.

For the maximal-order reference lattice
$U^2\perp\langle1,-3\rangle$, whose polarization product is $3$,
the exact arithmetic certificate gives $M=D$. Thus its fully
numerical cap is $D^6\mathcal B$.
The Gaussian sixfold extraction constant of HD-D03 gives
\begin{equation}
 c_t\le4634096\,\delta\,\lceil\sqrt T\rceil
                          M^{12}\mathcal B^2          \tag{R03.2}
\end{equation}
for a maximal-order target, with the same integral normalization
$(\delta,T)$ as HD-D01--D03. For a specified order-stable target
of conductor $f$, HD-S24 clears its actual lattice denominators
in $M$ and the integral extraction operator contributes the
additional factor $f^{24}$ in (R03.2).
This gives a terminating bound for each specified integral lattice;
no unstated height estimate replaces the computation of $M$.

\paragraph{Whole-proof checks.}
The root author read the complete argument, including the exact
trace composite, the complex group-algebra double centralizer,
the pure tensor evaluation, all degree costs and the actual
lattice transport. An independent complete review by
/root/imprimitive reconstructed each of these maps and bounds,
reopened Milne's full proofs, and recomputed the literal integer
relation and extraction ceilings. The full HD-G18 subgroup proof
has separate root and independent whole-proof checks, with all
107 distinct selected character realizations replayed. No theorem
from the candidate manuscript is a premise.
